% Solutions of 07-c-modules-make-tests.
\section*{Chapter \ref{ch:b1:c-modules-make-tests} --- Modular Programs, Makefiles and Tests}

\begin{solution}{exo:b1:c-modules-make-tests:1}
\code{check.h} is a \omterm{def:b1:c-modules-make-tests:makefile}{prerequisite} of \code{ratio_test.o} only: \code{make} compiles \code{ratio_test.c} and links
\code{ratio_test}, two commands. After editing \code{ratio.c}: compile it, then link both executables, three commands.
\end{solution}

\begin{solution}{exo:b1:c-modules-make-tests:2}
No: neither test makes \code{age < 18 || student} true without \code{age < 6}, so \code{price = 150} never runs. Adding
$(30, 1)$ (or any age from 6 to 17) runs it. The three tests then take both arcs of each decision: \code{age < 6} is
true for 4 and false for 70; the second condition is true for $(30, 1)$ and false for $(70, 0)$; \code{age >= 65} is
true for 70 and false for 4. \omterm{def:b1:c-modules-make-tests:coverage}{Branch coverage} is reached, counting each \code{if} as one decision.
\end{solution}

\begin{solution}{exo:b1:c-modules-make-tests:3}
The first call computes $100 \times 87 / 400 = 21$ in integer division, and \code{printf} sends \code{21} and a
percent sign to standard output; the second fails the \omterm{def:b1:c-modules-make-tests:assertion}{assertion}, which prints its message on standard error and stops
the program with exit status 134. As the output is a pipe here, the line \texttt{21\%} sat in the buffer of
\code{stdout} and was lost when \code{abort} stopped the program: only the message appears. With \code{-DNDEBUG},
\code{percent(22, 0)} divides by zero, which is \omterm{def:b1:c-types-and-control:ub}{undefined behaviour} (on this machine, the program is stopped by a
signal).
\omoutput{code/bachelor-1/07-c-modules-make-tests/outputs/exo_assert.txt}{The run.}
\end{solution}

\begin{solution}{exo:b1:c-modules-make-tests:4}
\code{exo_guard.c} includes \code{point.h}, then \code{segment.h}, which includes it again: \code{struct point} is
defined twice in the same file, an error. Put \texttt{\#ifndef POINT\_H}, \texttt{\#define POINT\_H} before the definition
and \texttt{\#endif} after it: the second inclusion then finds \code{POINT_H} defined and skips the text.
\end{solution}

\begin{solution}{exo:b1:c-modules-make-tests:5}
\code{limits_main.o} depends on \code{config.h}, which it includes, but the rule does not say so: \code{make} sees an
\omterm{def:b1:compilation-pipeline:compiler}{object file} newer than \code{limits_main.c} and does nothing. The rule must be \code{limits_main.o: limits_main.c
config.h}, or the dependencies written by the \omterm{def:b1:compilation-pipeline:compiler}{compiler} (\cref{exo:b1:c-modules-make-tests:11}).
\end{solution}

\begin{solution}{exo:b1:c-modules-make-tests:6}
The second test compares \code{c} with \code{a} instead of \code{m}: for \code{max3(1, 9, 5)}, \code{m} becomes 9, then
\code{5 > 1} replaces it by 5. The checks only say that the result is at least some argument, which a wrong result
can satisfy: \code{5 >= 5} passes. A test must compare with the exact expected value: \code{CHECK(max3(1, 9, 5) == 9)}
fails, as it should.
\omoutput{code/bachelor-1/07-c-modules-make-tests/outputs/exo_max3.txt}{Green, and wrong.}
\end{solution}

\begin{solution}{exo:b1:c-modules-make-tests:7}
\omcode{code/bachelor-1/07-c-modules-make-tests/exercises/stack.c}{5}{44}{A stack behind an \omterm{def:b1:c-modules-make-tests:opaque}{opaque type}}
The test \code{tests/stack_test.c} pops an empty stack, pushes 100 values (past the first capacities 4, 8, \dots),
pops them in reverse order, and frees a \omterm{def:b1:pointers-and-arrays:null}{null pointer}.
\end{solution}

\begin{solution}{exo:b1:c-modules-make-tests:8}
\code{(0, 1)}: the first decision true. \code{(3, 1)}: first false, \code{days_late > 30} false, cap false.
\code{(31, 2)}: both conditions of \code{&&} true, and the cap true (510 becomes 500). \code{(31, 7)}: the left
condition true, the right one false. These four take every arc, each atomic condition counted. The boundary:
\code{(30, 1)} gives 300 without surcharge, and \code{(31, 2)} is on the other side. The test is
\code{tests/late_fee_test.c}.
\end{solution}

\begin{solution}{exo:b1:c-modules-make-tests:9}
Every combination: $2^k$ tests. The $k + 1$ tests: all conditions true (outcome true), and for each $i$ the test where
only $c_i$ is false (outcome false). Each of the last $k$ differs from the first in $c_i$ alone, and the outcome
changes, so $c_i$ is shown to matter on its own. The test \code{tests/test_conditions.py} checks this for $k \leq 7$.
\end{solution}

\begin{solution}{exo:b1:c-modules-make-tests:10}
\code{make} rebuilds a \omterm{def:b1:c-modules-make-tests:makefile}{target} when a \omterm{def:b1:c-modules-make-tests:makefile}{prerequisite} is newer, and a \omterm{def:b1:c-modules-make-tests:makefile}{target} it rebuilds becomes newer than the \omterm{def:b1:c-modules-make-tests:makefile}{targets} that
depend on it: the rebuilding spreads along the arcs, from the edited file to every \omterm{def:b1:c-modules-make-tests:makefile}{target} reachable from it, and to no
other. \code{date.h}: \code{date.o}, \code{loans.o}, \code{loans_main.o}, \code{checker}; \code{date.c}: \code{date.o},
\code{checker}; \code{csv.h}: \code{csv.o}, \code{loans.o}, \code{checker}; \code{csv.c}: \code{csv.o}, \code{checker};
\code{loans.h}: \code{loans.o}, \code{loans_main.o}, \code{checker}; \code{loans.c} and \code{loans_main.c}: their \omterm{def:b1:compilation-pipeline:compiler}{object
file} and \code{checker}. The test \code{tests/test_conditions.py} computes these sets.
\end{solution}

\begin{solution}{exo:b1:c-modules-make-tests:11}
\omcode{code/bachelor-1/07-c-modules-make-tests/listings/mmd.Makefile}{2}{17}{Dependencies written by the \omterm{def:b1:compilation-pipeline:compiler}{compiler}}
\omoutput[6-14]{code/bachelor-1/07-c-modules-make-tests/outputs/exo_mmd.txt}{The fragment written for \code{loans.o}, and
the rebuild after touching \code{csv.h}.}
The \code{-include} line reads the \code{.d} files, and ignores them before the first build. The empty rules added by
\code{-MP} keep \code{make} from failing if a header is later deleted.
\end{solution}

\begin{solution}{exo:b1:c-modules-make-tests:12}
Take any statement. It lies on a \omterm{def:b1:c-modules-make-tests:cfg}{feasible path}; let $D$ be the last decision before it on that path, if any. If there
is none, every test executes it. Otherwise the statement lies on the arc of $D$ that the path takes, before the next
decision, and \omterm{def:b1:c-modules-make-tests:coverage}{branch coverage} means some test takes that arc, and so executes the statement. Converse: the function
\code{int f(int x) { int y = 0; if (x > 0) y = 1; return y; }} reaches \omterm{def:b1:c-modules-make-tests:coverage}{statement coverage} with the single test
$x = 1$, which never takes the false arc of the \code{if}.
\end{solution}

\begin{solution}{pb:b1:c-modules-make-tests:1}
\textbf{1.} \code{date.o} on \code{date.c}, \code{date.h}; \code{csv.o} on \code{csv.c}, \code{csv.h}; \code{loans.o}
on \code{loans.c}, \code{loans.h}, \code{date.h}, \code{csv.h}; \code{loans_main.o} on \code{loans_main.c},
\code{loans.h}, \code{date.h}; \code{checker} on the four \omterm{def:b1:compilation-pipeline:compiler}{object files}.

\textbf{2.} \code{struct loan} has fields of type \code{struct date}, whose size the \omterm{def:b1:compilation-pipeline:compiler}{compiler} must know. The \omterm{def:b1:c-modules-make-tests:interface}{interface}
of \code{loans} never mentions \code{csv}: splitting is a detail of its \omterm{def:b1:c-modules-make-tests:interface}{implementation}, so \code{csv.h} is included in
\code{loans.c} only, and users of \code{loans.h} do not depend on it.

\textbf{3.} It is a detail (another \omterm{def:b1:c-modules-make-tests:interface}{implementation} might count from another day), and keeping it private leaves the
module free to change it; its name cannot clash with a function of the same name elsewhere.

\textbf{4.} Opaque, a date would live on the heap, be created and freed by calls, and be read through functions: more
code, more allocations, and \code{struct loan} could hold only \omterm{def:b1:pointers-and-arrays:pointer}{pointers}. The gain: the representation could change (a
day count instead of three fields) without recompiling users. For three small integers, the cost wins.

\textbf{5.} \code{date} and \code{csv}: neither mentions loans.

\textbf{6.} ``Empty fields included'': a line ending with a comma yields a last field that is an empty string, so
\code{loan_parse} gets six fields and reads ``not returned'' from the empty sixth one.

\textbf{7.} Four \omterm{def:b1:compilation-pipeline:compiler}{compilations} (\code{-c}), then the link of \code{checker}.

\textbf{8.} \code{date.h} is included by \code{date.c}, \code{loans.c} (through \code{loans.h} as well) and
\code{loans_main.c} (through \code{loans.h}): 3. \code{csv.h}: \code{csv.o} and \code{loans.o}, 2. \code{loans.h}:
\code{loans.o} and \code{loans_main.o}, 2. Each \code{.c} file: its own \omterm{def:b1:compilation-pipeline:compiler}{object file}, 1.

\textbf{9.} $7 \times 4 = 28$ \omterm{def:b1:compilation-pipeline:compiler}{compilations} against $3 + 1 + 2 + 1 + 2 + 1 + 1 = 11$.

\textbf{10.} No: \code{loans_main.c} includes \code{loans.h}, which includes \code{date.h}, so a change to \code{date.h}
can change the meaning of \code{loans_main.c}. The \omterm{def:b1:c-modules-make-tests:makefile}{prerequisites} are the headers included directly or indirectly.

\textbf{11.} An edit to \code{date.h} that changes \code{struct date} (a field added, say) would leave \code{loans.o}
compiled with the old layout while \code{date.o} uses the new one: the two would read the same bytes as different
fields, and the program would compute wrong delays without any error.

\textbf{12.} Nothing: the counts concern the \omterm{def:b1:c-modules-make-tests:makefile}{target} \code{checker}. With \code{date_test} in the default \omterm{def:b1:c-modules-make-tests:makefile}{target}, an edit
to \code{date.h} or \code{date.c} would add one \omterm{def:b1:compilation-pipeline:compiler}{compilation} (\code{date_test.o} for the first) and one link.

\textbf{13.} Years not divisible by 4 (2025); divisible by 4 but not by 100 (2024); divisible by 100 but not by 400
(1900, 2100); divisible by 400 (2000, 2400).

\textbf{14.} \code{2025-02-29} and \code{2024-02-29} (the last day of February, each side of the leap rule);
\code{2025-13-01} and \code{2025-00-10} (the months 1 and 12); \code{2025-04-31} and \code{2025-04-00} (the days 1 and 30
of a 30-day month); \code{1969-12-31} (the start of the domain); \code{2025-12-31} to \code{2026-01-01} (a year's end).

\textbf{15.} From 1970 to 1999: 30 years, of which 1972, 1976, \dots, 1996 are leap years, 7 of them;
$30 \times 365 + 7 = 10\,957$.

\textbf{16.} The two arcs of the \omterm{def:b1:c-modules-make-tests:assertion}{assertion} in \code{days_in_month} that lead to its failure (\code{month >= 1} false,
\code{month <= 12} false). Every caller passes a valid month: \code{date_parse} checks it first, and the loop of
\code{days_since_1970} runs from 1 to 11. A failing \omterm{def:b1:c-modules-make-tests:assertion}{assertion} would mean a bug in the module.

\textbf{17.} None of them: the suite tests the date module only. The figures are produced by the whole program on the
real file, and only a test of \code{loans} and \code{loans_main} (or a comparison with a count made another way) would
check them.

\textbf{18.} \code{loan_days_late} might use \code{out} instead of \code{due}, or \code{loans_main.c} might skip the
last line of the file: the dates would be right and the total wrong. \omterm{def:b1:c-modules-make-tests:unit-test}{Unit tests} of one module say nothing of the
modules around it.

\textbf{19.} \code{is_leap} returns \texttt{year \% 4 == 0}, which makes 1900 and 2100 leap years. The check of line 17
calls \code{is_leap(2100)} directly; the check of line 34 counts 367 days from 2099-02-28 to 2100-03-01, because the
first version gives February 2100 a 29th day.

\textbf{20.} No: between 2025 and 2026 the two rules agree, every year from 1901 to 2099 being leap exactly when
divisible by 4. The bug was found because the tests were chosen from the classes of the specification, not from the
data: the class ``divisible by 100 but not by 400'' had a test.

\textbf{21.} A month 13 can come from outside the program (a typing error in the file): \code{date_parse} must expect
it and return an error. Once a date is parsed, a month 13 can only come from a bug: \code{days_in_month} asserts.

\textbf{22.} A check: the number comes from the file, written by people and other programs. The checker reports the line
and stops with an error status.

\textbf{23.} \code{date_parse} rejects 2025-02-29, so \code{main} prints its usage line on standard error and returns 2.

\textbf{24.} \emph{The rebuild saving.} With the four modules and the dependencies of
\cref{lst:b1:c-modules-make-tests:library-mk}, one edit to each of the seven \omterm{def:b1:compilation-pipeline:source-file}{source files} costs $3 + 1 + 2 + 1 + 2 + 1 +
1 = 11$ \omterm{def:b1:compilation-pipeline:compiler}{compiler} runs instead of $7 \times 4 = 28$ for full rebuilds: an edit to a file costs as many \omterm{def:b1:compilation-pipeline:compiler}{compilations} as
there are \omterm{def:b1:compilation-pipeline:compiler}{object files} reachable from it in the dependency graph, and \code{make} counted exactly those.
\end{solution}
